% !TEX program = xelatex
% Lernplatt - by Can Siebert
% Kurs 71 (Dig 42) - Tag 4 - Lehrskript
% Repository, Governance, Kollaboration; Execution-Modelle (Automatisierung, Workflows)
%
% Self-contained xelatex source. Kein biblatex, keine externen Bilder.

\documentclass[11pt,a4paper,oneside]{article}

% --- Encoding & Sprache --------------------------------------------------
\usepackage{polyglossia}
\setdefaultlanguage{german}

% --- Geometrie -----------------------------------------------------------
\usepackage[a4paper,top=2cm,bottom=2cm,outer=2.5cm,inner=2.5cm,bindingoffset=1.5cm]{geometry}

% --- Fonts ---------------------------------------------------------------
\usepackage{fontspec}
\defaultfontfeatures{Ligatures=TeX}

% Charter als Body. Fallback: TeX Gyre Schola (sehr aehnlich, in MiKTeX dabei).
\IfFontExistsTF{Charter}{%
  \setmainfont{Charter}[Numbers=OldStyle]%
}{%
  \IfFontExistsTF{Noto Serif}{%
    \setmainfont{Noto Serif}%
  }{%
    \setmainfont{Liberation Serif}%
  }%
}

% Sans = Source Sans 3 / Source Sans Pro. Fallback: TeX Gyre Heros.
\IfFontExistsTF{Source Sans 3}{%
  \setsansfont{Source Sans 3}%
}{%
  \IfFontExistsTF{Source Sans Pro}{%
    \setsansfont{Source Sans Pro}%
  }{%
    \setsansfont{Liberation Sans}%
  }%
}

% Caveat fuer Akzente. Wenn nicht da -> faellt auf italic Hauptschrift zurueck.
\newcommand{\caveatfontfallback}{\itshape}
\IfFontExistsTF{Caveat}{%
  \newfontfamily\caveatfont{Caveat}[Scale=1.15]%
}{%
  \newcommand\caveatfont{\caveatfontfallback}%
}

% Monospace
\setmonofont{Liberation Mono}

% --- Mikro-Typografie ----------------------------------------------------
\usepackage{microtype}
\usepackage{eurosym}
\linespread{1.15}

% --- Farben & Boxen ------------------------------------------------------
\usepackage{xcolor}
\definecolor{lpInk}{RGB}{20,28,38}
% Kurs 71 Akzent: Orange (#9d4a1a)
\definecolor{kurs71}{HTML}{9D4A1A}
\definecolor{lpAccent}{HTML}{9D4A1A}
\definecolor{lpAccentLight}{RGB}{248,232,218}
\definecolor{lpAmber}{RGB}{222,138,30}
\definecolor{lpAmberLight}{RGB}{255,243,217}
\definecolor{lpAmberBorder}{RGB}{196,118,18}
\definecolor{lpGray}{RGB}{96,104,114}
\definecolor{lpGrayLight}{RGB}{238,240,243}

\usepackage[most]{tcolorbox}
\tcbuselibrary{breakable,skins}

% Eselsbruecke - amber, dashed
\newtcolorbox{eselsbox}[1][Eselsbrücke]{%
  enhanced, breakable,
  colback=lpAmberLight,
  colframe=lpAmberBorder,
  boxrule=0.6pt,
  arc=2pt,
  borderline={0.4pt}{0pt}{lpAmberBorder, dashed},
  left=10pt,right=10pt,top=8pt,bottom=8pt,
  fonttitle=\sffamily\bfseries\color{lpAmberBorder},
  coltitle=lpAmberBorder,
  title={#1},
  attach boxed title to top left={xshift=8pt,yshift=-8pt},
  boxed title style={size=small, colback=lpAmberLight, colframe=lpAmberBorder, boxrule=0.4pt, arc=1pt},
  top=14pt
}

% Merkkasten - orange fuer Kurs 71
\newtcolorbox{merkbox}[1][Merksatz]{%
  enhanced, breakable,
  colback=lpAccentLight,
  colframe=lpAccent,
  boxrule=0.6pt,
  arc=2pt,
  left=10pt,right=10pt,top=8pt,bottom=8pt,
  fonttitle=\sffamily\bfseries\color{white},
  coltitle=white,
  title={#1},
  attach boxed title to top left={xshift=8pt,yshift=-8pt},
  boxed title style={size=small, colback=lpAccent, colframe=lpAccent, boxrule=0.4pt, arc=1pt},
  top=14pt
}

% --- Listen --------------------------------------------------------------
\usepackage{enumitem}
\setlist{itemsep=2pt, topsep=4pt, parsep=0pt}
\setlist[itemize,1]{label={\textbullet},leftmargin=1.6em}
\setlist[itemize,2]{label={\textendash},leftmargin=1.4em}
\setlist[enumerate,1]{label=\arabic*., leftmargin=1.8em}

% --- Tabellen ------------------------------------------------------------
\usepackage{tabularx}
\usepackage{array}
\usepackage{booktabs}
\usepackage{longtable}

% --- Kopf-/Fusszeilen ----------------------------------------------------
\usepackage{fancyhdr}
\usepackage{lastpage}
\pagestyle{fancy}
\fancyhf{}
\renewcommand{\headrulewidth}{0.3pt}
\renewcommand{\footrulewidth}{0pt}
\fancyhead[L]{\footnotesize\sffamily\color{lpGray}Lernplatt \textperiodcentered{} Kurs 71 \textperiodcentered{} Tag 4}
\fancyhead[R]{\footnotesize\sffamily\color{lpGray}Repository \textperiodcentered{} Governance \textperiodcentered{} Execution}
\fancyfoot[L]{\footnotesize\sffamily\color{lpGray}\textcopyright{} Can Siebert}
\fancyfoot[R]{\footnotesize\sffamily\color{lpGray}\thepage{} / \pageref{LastPage}}

\fancypagestyle{plain}{%
  \fancyhf{}%
  \renewcommand{\headrulewidth}{0pt}%
  \fancyfoot[C]{\footnotesize\sffamily\color{lpGray}\thepage{} / \pageref{LastPage}}%
}

% --- Section-Styling -----------------------------------------------------
\usepackage[explicit]{titlesec}

\titleformat{\section}[block]
  {\sffamily\Large\bfseries\color{lpAccent}}
  {\thesection.}{0.6em}{#1}
  [{\vspace{2pt}\color{lpAccent}\titlerule[0.6pt]}]

\titleformat{\subsection}[block]
  {\sffamily\large\bfseries\color{lpInk}}
  {\thesubsection}{0.6em}{#1}

\titleformat{\subsubsection}[block]
  {\sffamily\normalsize\bfseries\color{lpInk}}
  {}{0pt}{#1}

\titlespacing*{\section}{0pt}{14pt}{8pt}
\titlespacing*{\subsection}{0pt}{10pt}{4pt}
\titlespacing*{\subsubsection}{0pt}{8pt}{2pt}

% --- Hyperlinks ----------------------------------------------------------
\usepackage[hidelinks,unicode]{hyperref}

% --- TikZ ----------------------------------------------------------------
\usepackage{tikz}
\usetikzlibrary{shapes.geometric,arrows.meta,positioning,calc,fit,backgrounds}
\definecolor{lpGreenLight}{RGB}{225,238,225}

% --- TOC -----------------------------------------------------------------
\renewcommand{\contentsname}{\sffamily\Large\bfseries\color{lpAccent}Inhalt}

% --- Helfer --------------------------------------------------------------
\newcommand{\akzent}[1]{{\caveatfont\color{lpAmber}#1}}
\newcommand{\fachbegriff}[1]{\textbf{#1}}
\newcommand{\quelle}[1]{{\footnotesize\color{lpGray}(#1)}}

% =========================================================================
\begin{document}

% --- COVER ---------------------------------------------------------------
\begin{titlepage}
  \pagestyle{empty}
  \vspace*{1.5cm}
  \noindent{\sffamily\footnotesize\color{lpGray}LERNPLATT \textperiodcentered{} BY CAN SIEBERT}\\[2pt]
  \noindent{\color{lpAccent}\rule{\linewidth}{1.2pt}}\\[4pt]
  \noindent{\sffamily\footnotesize\color{lpGray}MODUL 2 \textperiodcentered{} TAG 4 VON 4 \textperiodcentered{} KURS 71 (Dig 42) \textperiodcentered{} 2.\,Juni 2026}

  \vspace{3cm}
  {\sffamily\fontsize{30}{34}\selectfont\bfseries\color{lpInk}Repository\\\&\\Execution\par}

  \vspace{0.7cm}
  {\sffamily\large\color{lpGray}Vom Prozessbild zum steuerbaren System --\\Single Source of Truth, Governance und ausführbare Workflows.\par}

  \vspace{1.5cm}
  {\caveatfont\LARGE\color{lpAmber}Lehrskript -- Tag 4 von 16}\par

  \vfill

  \noindent\begin{minipage}{0.55\linewidth}
    {\sffamily\footnotesize\color{lpGray}Lernpfad:}\\
    {\sffamily\small\color{lpInk}Wissen \textrightarrow{} Anwendung \textrightarrow{} Management}\\[10pt]
    {\sffamily\footnotesize\color{lpGray}Bearbeitungsdauer:}\\
    {\sffamily\small\color{lpInk}1 Kurstag}
  \end{minipage}\hfill
  \begin{minipage}{0.4\linewidth}
    \raggedleft
    {\sffamily\footnotesize\color{lpGray}Dozent:}\\
    {\sffamily\small\color{lpInk}Can Siebert}\\[10pt]
    {\sffamily\footnotesize\color{lpGray}Format:}\\
    {\sffamily\small\color{lpInk}Theorie \textperiodcentered{} Fallstudie \textperiodcentered{} Coaching}
  \end{minipage}

  \vspace{0.5cm}
  \noindent{\color{lpAccent}\rule{\linewidth}{0.6pt}}
\end{titlepage}

% --- LERNZIELE -----------------------------------------------------------
\section*{Lernziele dieses Tages}
\addcontentsline{toc}{section}{Lernziele dieses Tages}
\thispagestyle{plain}

\noindent An den ersten drei Tagen haben Sie Prozesse aus mehreren Perspektiven \emph{modelliert} -- als Notation, als Wertstrom, als Datenrealität, mit Blick auf Stakeholder, KPIs und SOA. Heute machen wir den entscheidenden Schritt von der einzelnen Diagrammdatei zum \emph{betriebenen System}: einem Repository, das Modelle steuerbar macht, und einer Execution-Schicht, die Prozesse ausführbar und auditierbar werden lässt.

\subsection*{L1 -- Wissen (Reproduktion und Einordnung)}
\begin{itemize}
  \item Den Repository-Ansatz mit seinen Bestandteilen (Prozesslandkarte, Modelle, Rollen, Dokumente, Versionen, KPIs, Risiken) einordnen.
  \item Modell-Governance benennen: Owner, Modeler, Reviewer, Approver -- und das Zusammenspiel mit Benennungs-, Versions- und Freigaberegeln.
  \item Die Differenz zwischen \emph{Business-BPMN} (verständlich, kommuniziert Verantwortung) und \emph{Execution-BPMN} (präzise, maschinen-nah) erklären.
  \item Camunda-typische Konstrukte einordnen: User Task, Service Task, Timer-/Message-/Error-Event, Datenobjekte und Variablen.
\end{itemize}

\subsection*{L2 -- Anwendung (Transfer auf konkrete Situationen)}
\begin{itemize}
  \item Eine Repository-Struktur in drei Ebenen (E2E-Prozess \textrightarrow{} Subprozess \textrightarrow{} Arbeitsdokument) für einen Beispielprozess entwerfen.
  \item Einen Mini-Prozess in einen \emph{Automations-Slice} schneiden: Start, Ende, Datenminimum, Ausnahmepolitik definieren.
  \item Variantenführung bewusst entscheiden: separate Modelle oder Variantenattribute -- mit Trade-off-Begründung.
  \item KPIs für ein Workflow-Programm definieren (Durchlaufzeit, Eskalationsquote, Rework) und Owner zuordnen.
\end{itemize}

\subsection*{L3 -- Management (Entscheidung und Verantwortung)}
\begin{itemize}
  \item Repository als \emph{Betriebssystem} für Prozesssteuerung verstehen -- nicht als Ablage.
  \item Execution als \emph{Vertrag} zwischen Fach, IT und Compliance behandeln: Daten, Fehler, SLAs (Service Level Agreements), Eskalation.
  \item Entscheidungen unter Unsicherheit treffen: MVP (\emph{Minimum Viable Product}, erste marktreife Mindestversion) statt Perfektion, mit Frühwarnsignalen abgesichert.
  \item Governance-Regeln auf das Minimum reduzieren, das 80\,\% des Chaos verhindert -- und alles weitere bewusst delegieren.
\end{itemize}

\begin{merkbox}[Roter Faden]
Modelle leben oder sterben \emph{nach} der Modellierung -- in einem Repository mit Governance, oder ausgeführt in einer Engine mit Audit Trail. Wer modelliert, ohne den Lebenszyklus mitzudenken, baut Diagramme für Schubladen. Senior-Praxis heißt: \emph{Repository steuert Transparenz, Execution steuert Durchlaufzeit -- und beide brauchen denselben Owner.}
\end{merkbox}

\clearpage

% --- 1. EINSTIEG ---------------------------------------------------------
\section{Einstieg: Vom Diagramm zum steuerbaren System}

Es gibt eine bemerkenswerte Konstante in fast jeder Organisation, die seit ein paar Jahren modelliert hat: Es existieren \emph{viele} Diagramme. Sie liegen in SharePoint-Ordnern, in Visio-Dateien, in Powerpoint-Anhängen, in PDF-Exporten. Was meist fehlt, ist die Fähigkeit, sie als \emph{System} zu betreiben. Niemand weiß, welches die aktuelle Version ist; niemand weiß, wer entscheidet, ob eine Änderung freigegeben wird; niemand weiß, ob das, was im Diagramm steht, mit dem übereinstimmt, was die Workflow-Engine ausführt. \quelle{Dumas et al., 2018, Kap.~1; Harmon, 2019}

Der heutige Tag adressiert genau diese Lücke. Modellieren ist eine notwendige Bedingung, aber keine hinreichende, um Prozesse zu \emph{steuern}. Hinreichend wird es erst durch zwei Schichten, die typischerweise unterschätzt werden: eine \fachbegriff{Repository-Schicht}, die Modelle als verwaltete Artefakte hält, und eine \fachbegriff{Execution-Schicht}, die ausgewählte Modelle tatsächlich ausführbar macht.

\subsection{Drei Symptome eines fehlenden Repository}

\begin{enumerate}
  \item \fachbegriff{Niemand findet die aktuelle Version.} Ein Auditor fragt, welcher Stand verbindlich ist. Drei Antworten von drei Personen, alle leicht unterschiedlich.
  \item \fachbegriff{Standorte modellieren divergent.} Jede Niederlassung benennt anders, modelliert in unterschiedlichen Tiefen, dokumentiert auf eigenen Laufwerken. Vergleichbarkeit ist Glücksspiel.
  \item \fachbegriff{Änderungen passieren ``irgendwie''.} Eine Anpassung des Genehmigungsworkflows landet als E-Mail-Notiz im Postfach, nicht als versionierter Change Request mit Reviewer-Verantwortung.
\end{enumerate}

Diese Symptome sind nicht primär Tool-Probleme. Sie sind \emph{Governance-Probleme}, die durch ein gutes Tool unterstützt -- aber nicht ersetzt -- werden können.

\subsection{Was ein Repository leisten muss}

Ein \fachbegriff{Prozess-Repository} ist mehr als ein Ordner mit Diagrammen. Es ist eine kuratierte Sammlung von Artefakten mit definierten Metadaten, Verantwortlichkeiten und Lebenszyklen. Im Kern braucht es: \quelle{Scheer, 2002; Rosemann \& vom Brocke, 2010}

\begin{itemize}
  \item Eine \fachbegriff{Prozesslandkarte} (Value Chain), die die End-to-End-Sicht ordnet -- typischerweise 8 bis 15 Kernprozesse auf der obersten Ebene.
  \item Die zugehörigen \fachbegriff{Prozessmodelle} (BPMN, EPC, Activity), Subprozesse und Arbeitsdokumente.
  \item \fachbegriff{Rollen}: Process Owner, Modeler, Reviewer, Approver. Pflicht, nicht optional.
  \item \fachbegriff{Metadaten}: Owner, Gültigkeit, Status (Draft, Review, Approved, Retired), Version, letzte Änderung, Standort, Risikoklasse.
  \item Optional: Verlinkung zu \fachbegriff{KPIs}, \fachbegriff{Risiken}, \fachbegriff{Controls} -- damit das Repository nicht nur Diagramme hält, sondern \emph{Steuerungsobjekte}.
\end{itemize}

\begin{eselsbox}[Eselsbrücke: \akzent{P-M-R-V}]
Ein Repository steht auf vier Beinen, das fünfte ist der Owner:\\[2pt]
\textbf{P}rozesslandkarte -- die Übersicht über alle Kernprozesse.\\
\textbf{M}odelle -- BPMN, EPC, Subprozesse, Arbeitsdokumente.\\
\textbf{R}ollen -- Owner, Modeler, Reviewer, Approver.\\
\textbf{V}ersionen -- Status, Gültigkeit, Änderungslog.\\[2pt]
\emph{Wer ein Bein vergisst, steht auf einem Tisch, der wackelt. Wer alle fünf hat, hat ein Betriebssystem für Prozesse.}
\end{eselsbox}

\subsection{Was Execution leisten muss}

Wenn das Repository der ``Bibliothekar'' ist, ist die Execution-Schicht die ``Maschine''. Sie nimmt ausgewählte Modelle und führt sie aus -- als Workflows mit User Tasks, Service Tasks, Timern, Eskalationen und Audit Trail. \quelle{Freund \& Rücker, 2019; OMG, 2014}

Nicht jedes Modell muss ausführbar sein. Im Gegenteil: Die wertvollsten Modelle sind oft die \emph{Übersichtsmodelle}, die Verantwortung sichtbar machen und Kommunikation strukturieren. Aber bestimmte Teilstrecken -- typischerweise solche mit klarer Start-Ende-Definition, eindeutigen Daten und nachvollziehbaren Fehlerfällen -- profitieren massiv davon, in einer \fachbegriff{Workflow-Engine} (Camunda, Flowable, jBPM, oder vergleichbare BPMS-Produkte; BPMS = \emph{Business Process Management System}) zu laufen.

% --- 2. REPOSITORY-DESIGN ------------------------------------------------
\section{Repository-Design: Single Source of Truth in drei Ebenen}

Ein praxistaugliches Repository kommt mit einer überraschend einfachen Struktur aus. Drei Ebenen reichen für den Start -- mehr Ebenen erzeugen Pflegeaufwand, der den Nutzen oft übersteigt.

\subsection{Die drei kanonischen Ebenen}

\begin{enumerate}
  \item \fachbegriff{Ebene 1 -- E2E-Prozess.} Die Sicht auf einen vollständigen Wertschöpfungsstrang: ``Order-to-Cash'', ``Incident-to-Resolution'', ``Lead-to-Order'', ``Hire-to-Retire''. Auf dieser Ebene wird \emph{nicht} im Detail modelliert, sondern Scope, Owner, KPI (\emph{Key Performance Indicator}) und übergeordneter Verlauf geklärt.
  \item \fachbegriff{Ebene 2 -- Subprozess.} Die Detail-Modelle, die zeigen, wie ein Teil des E2E-Prozesses tatsächlich abläuft. Hier wohnt das eigentliche BPMN/EPC-Diagramm mit Lanes, Gateways und Tasks. Pro E2E-Prozess typischerweise 3 bis 8 Subprozesse.
  \item \fachbegriff{Ebene 3 -- Arbeits- und Nachweisdokument.} Die operativen Artefakte: SOP (Standard Operating Procedure), Checklisten, Formulare, Anweisungen, Schulungsunterlagen, Auditnachweise. Sie hängen an einzelnen Subprozessen oder Tasks.
\end{enumerate}

\subsection{Die Mindestmenge an Metadaten}

Wer Modelle vergleichbar halten will, muss Pflichtfelder definieren. \emph{Fünf} Felder sind ein guter Anfang -- mehr ist Komfort, weniger ist Schluderei: \quelle{Allweyer, 2020; Harmon, 2019}

\begin{itemize}
  \item \textbf{Owner.} Genau eine Person (Rolle), nicht ein Verteiler. ``Wer eskaliert, ruft Owner an.''
  \item \textbf{Gültigkeit.} Ab wann gilt diese Version, bis wann ist sie aktuell, wann ist Review fällig?
  \item \textbf{Status.} Draft, Review, Approved, Retired. Keine ``halben'' Zustände dazwischen.
  \item \textbf{Standort/Variante.} Welche Organisation, welche Region, welche Produkt-Variante?
  \item \textbf{Risikoklasse.} Hoch (compliance- oder geldkritisch), Mittel, Niedrig -- steuert Review-Tiefe.
\end{itemize}

\begin{merkbox}[Owner ist eine Person, nicht eine Abteilung]
Ein Process Owner ist eine \emph{konkrete Rolle} mit Namen, nicht ``Bereich Vertrieb''. Wer als Owner einen Verteiler einträgt, hat de facto keinen Owner. Im Audit-Fall stellt sich heraus: niemand fühlt sich zuständig, weil alle gemeint waren und keiner verantwortet hat.
\end{merkbox}

\subsection{Variantenführung -- als Attribut oder als eigenes Modell?}

Die schwierigste Repository-Entscheidung ist meist nicht die Struktur, sondern die \emph{Variantenpolitik}. Drei Standorte machen das Onboarding leicht unterschiedlich; zwei Produktlinien folgen ähnlichen, aber nicht identischen Prozessen. Sollen das separate Modelle sein, oder ein Modell mit Variantenattributen?

Faustregel:
\begin{itemize}
  \item \textbf{Variantenattribut} -- wenn 80\,\% des Modells identisch ist und nur Detail-Felder (Schwellwert, Approver-Rolle, SLA) variieren. Vorteil: Pflege einmal, Konsistenz hoch. Nachteil: Modell wirkt überladen.
  \item \textbf{Separates Modell} -- wenn Hauptpfade abweichen, andere Akteure beteiligt sind oder regulatorische Unterschiede bestehen. Vorteil: Klarheit. Nachteil: Pflegeaufwand multipliziert sich.
\end{itemize}

\subsection{Repository-Blueprint -- ein Beispiel}

Für einen E2E-Prozess ``Order-to-Cash'' kann ein realistisches Repository so aussehen:

\begin{tabularx}{\linewidth}{@{}>{\bfseries}lX@{}}
\toprule
Ebene 1 & Order-to-Cash (Owner: Head of Sales Ops, KPI: DLZ 7\,Tage, Status: Approved v3.1) \\
Ebene 2 & Auftragsannahme, Verfügbarkeitsprüfung, Kreditcheck, Auslieferung, Fakturierung, Mahnwesen \\
Ebene 3 & SOP Verfügbarkeitsprüfung, Eskalationsleitfaden Kreditcheck, Mahnstufen-Brieftexte, Audit-Nachweis Fakturierung \\
\bottomrule
\end{tabularx}

\medskip

\noindent Diese Sicht ist nicht spektakulär -- aber sie ist \emph{vollständig}. Sie macht sichtbar, wer im Notfall anzurufen ist, welche Version gilt, und wo das nächste Detail liegt. Das ist der Unterschied zwischen einer Ordnerstruktur und einem Steuerungsinstrument.

% --- 3. GOVERNANCE & ROLLEN ----------------------------------------------
\section{Governance: Rollen, Reviews und Freigaben}

Ohne Rollen wird ein Repository binnen Monaten zur Abstellkammer. \emph{Mit} Rollen, klaren Regeln und einem ehrlichen Review wird es zum Steuerungssystem. Die Verteilung darf dabei pragmatisch sein: vier Rollen sind in der Praxis ausreichend, mehr macht es schwerfällig. \quelle{Rosemann \& vom Brocke, 2010}

\subsection{Die vier Kernrollen}

\begin{itemize}
  \item \fachbegriff{Process Owner.} Verantwortet den E2E-Prozess fachlich, vertritt ihn nach außen, entscheidet über strategische Änderungen. Hängt an Ebene 1.
  \item \fachbegriff{Modeler.} Modelliert konkret, in Abstimmung mit den Fachbereichen. Beherrscht das Handwerk. Hängt an Ebene 2.
  \item \fachbegriff{Reviewer.} Prüft Modelle gegen Konventionen, GoM und 7PMG. Mindestens vier Augen pro Approval.
  \item \fachbegriff{Approver.} Gibt finale Modellversionen frei. Häufig der Process Owner selbst, oder ein Bereichsleiter mit fachlicher Verantwortung.
\end{itemize}

\subsection{Der Freigabeprozess in drei Stufen}

Ein einfacher und in jeder Organisation tragfähiger Lebenszyklus:

\begin{enumerate}
  \item \fachbegriff{Draft.} Modeler arbeitet am Modell, schreibt Annahmen sichtbar in Annotations.
  \item \fachbegriff{Review.} Reviewer prüft auf Konsistenz, Vollständigkeit, Konventionen. Feedback zurück an Modeler, ggf. mehrere Iterationen.
  \item \fachbegriff{Approved.} Approver gibt frei. Version wird mit Datum, Reviewer-Namen und Approver-Namen versiegelt. Nur diese Version ist verbindlich für Operation und Audit.
\end{enumerate}

Hinzu kommt ein vierter Status: \fachbegriff{Retired}, wenn ein Modell durch eine neuere Version oder einen veränderten Prozess obsolet wird. Wichtig: \emph{Retired-Modelle nicht löschen}, sondern archivieren -- für Audit-Rückfragen.

\begin{eselsbox}[Eselsbrücke: \akzent{D-R-A-R}]
Vier Stationen im Modell-Lebenszyklus, in dieser Reihenfolge:\\[2pt]
\textbf{D}raft -- Modeler arbeitet, Annahmen sichtbar.\\
\textbf{R}eview -- Reviewer prüft gegen Konvention.\\
\textbf{A}pproved -- Approver versiegelt mit Datum und Namen.\\
\textbf{R}etired -- archivieren, nicht löschen.\\[2pt]
\emph{Wer eine Station überspringt, sammelt Schulden im Repository -- die spätestens im Audit fällig werden.}
\end{eselsbox}

\subsection{Modellierungs-Policy auf einer Seite}

Eine effektive \fachbegriff{Modellierungs-Policy} passt auf eine Seite und regelt das Minimum. Mehr ist Bürokratie:

\begin{enumerate}
  \item Notation: BPMN 2.0 für ausführungsnahe Modelle, EPC oder Activity für Übersichten.
  \item Benennung: Verb + Objekt für Aktivitäten/Funktionen, Zustand für Ereignisse, ``X-to-Y'' für E2E-Prozesse.
  \item Pflicht-Metadaten: Owner, Status, Version, Gültigkeit, Risikoklasse.
  \item Pflicht-Elemente in jedem Modell: Start- und End-Ereignis, mindestens eine Lane mit klarer Rolle, jede Gateway-Bedingung beschriftet.
  \item Review-Pflicht: vor jedem Approved-Status, ohne Ausnahme.
  \item Variantenregel: Attribute bevorzugt, separate Modelle nur bei begründeter Hauptpfad-Divergenz.
  \item Update-Rhythmus: jedes Modell mindestens jährlich reviewen, bei Prozess-Änderungen sofort.
\end{enumerate}

\subsection{Die zwei Regeln, die 80\,\% des Chaos verhindern}

Wenn man nur zwei Governance-Regeln durchsetzen darf, sind es diese: \quelle{Harmon, 2019}

\begin{itemize}
  \item \emph{Jedes Modell hat einen genannten Owner und einen Approved-Status mit Datum.} Das eliminiert Versions-Verwirrung und Verantwortungs-Vakuum.
  \item \emph{Änderungen laufen über einen Change Request mit Reviewer.} Das eliminiert die ``mal eben angepasst''-Kategorie, die Audit-Findings produziert.
\end{itemize}

Alles weitere ist Komfort und kann später ergänzt werden, wenn die Organisation gelernt hat, mit dem Minimum zu leben.

% --- 4. BUSINESS- VS. EXECUTION-BPMN -------------------------------------
\section{Business-BPMN vs.\ Execution-BPMN}

Eine der häufigsten Quellen für Frust in Modellierungs-Projekten ist die unsichtbare Annahme, dass \emph{ein} Modell gleichzeitig dem Vorstand etwas erklären und der Workflow-Engine etwas vorgeben kann. In der Praxis sind das zwei verschiedene Modelle, die denselben Prozess beschreiben -- mit unterschiedlicher Tiefe und Strenge. \quelle{Freund \& Rücker, 2019; OMG, 2014}

\begin{figure}[h]
\centering
\begin{tikzpicture}[font=\sffamily\footnotesize,
  ev/.style={circle, draw=lpGray, line width=0.6pt, minimum size=4mm, inner sep=0pt},
  evEnd/.style={circle, draw=lpGray, line width=1.2pt, minimum size=4mm, inner sep=0pt},
  task/.style={rectangle, rounded corners=2pt, draw=lpAccent, fill=lpAccentLight,
    line width=0.5pt, minimum width=14mm, minimum height=7mm, align=center, inner sep=1pt},
  gw/.style={diamond, draw=lpAccent, fill=white, line width=0.5pt,
    minimum size=5mm, inner sep=0pt},
  arr/.style={-{Latex[length=1.8mm]}, draw=lpGray, line width=0.5pt},
  lbl/.style={font=\sffamily\scriptsize\bfseries, color=lpAccent, align=center},
  hint/.style={font=\sffamily\scriptsize\itshape, color=lpGray, align=center}]
  % Business
  \node[lbl] at (3.7,1.5) {Business-BPMN\\\scriptsize\mdseries\upshape Verständigung};
  \node[ev] (b0) at (0,0) {};
  \node[task] (b1) at (1.4,0) {Anfrage\\prüfen};
  \node[gw] (b2) at (3.5,0) {};
  \node[task] (b3) at (5.5,0) {bearbeiten};
  \node[evEnd] (b4) at (7.3,0) {};
  \draw[arr] (b0) -- (b1); \draw[arr] (b1) -- (b2); \draw[arr] (b2) -- (b3); \draw[arr] (b3) -- (b4);
  % Execution (Slice)
  \node[lbl] at (3.7,-1.6) {Execution-BPMN (Slice)\\\scriptsize\mdseries\upshape Ausführung};
  \node[ev] (e0) at (1.4,-2.7) {};
  \node[task, fill=lpGreenLight, draw=lpGray] (e1) at (3.0,-2.7) {Service\\Task (API)};
  \node[task, fill=lpGreenLight, draw=lpGray] (e2) at (5.0,-2.7) {User\\Task};
  \node[evEnd] (e3) at (6.6,-2.7) {};
  \draw[arr] (e0) -- (e1); \draw[arr] (e1) -- (e2); \draw[arr] (e2) -- (e3);
  % Slice bracket
  \draw[lpAccent, dashed, line width=0.5pt] (3.0,-0.5) -- (3.0,-2.2);
  \draw[lpAccent, dashed, line width=0.5pt] (5.0,-0.5) -- (5.0,-2.2);
  \node[hint] at (4.0,-1.0) {Automations-Slice};
\end{tikzpicture}
\caption{\sffamily Vom Business-BPMN zum ausführbaren Slice.}
\label{fig:konzept-71t04}
\end{figure}

\subsection{Business-BPMN: verständlich, kommuniziert Verantwortung}

Ein \fachbegriff{Business-BPMN}-Modell ist für die \emph{Fachseite} optimiert. Es darf an manchen Stellen unscharf sein, solange es:

\begin{itemize}
  \item den Hauptpfad sauber zeigt,
  \item Verantwortlichkeiten über Lanes sichtbar macht,
  \item Entscheidungspunkte (Gateways) eindeutig benennt,
  \item Schnittstellen zu anderen Pools markiert,
  \item Annahmen explizit dokumentiert.
\end{itemize}

Ziel ist Verständigung -- nicht Vollständigkeit. Ein Business-BPMN kann an einzelnen Stellen Sub-Prozesse nutzen, ohne sie zu detaillieren, oder Manual Tasks beschreiben, ohne sie ausführbar zu machen.

\subsection{Execution-BPMN: präzise, maschinen-nah}

Ein \fachbegriff{Execution-BPMN}-Modell muss in eine Workflow-Engine laden und dort \emph{ausgeführt} werden können. Das verlangt eine andere Disziplin: jede Bedingung formal, jeder Pfad vollständig, jede Variable typisiert. Konkret kommen folgende Elemente in den Vordergrund:

\begin{itemize}
  \item \fachbegriff{User Task} -- ein menschlicher Schritt mit Form, Zuweisungsregel und SLA.
  \item \fachbegriff{Service Task} -- ein Systemschritt, typischerweise ein API-Aufruf (API = \emph{Application Programming Interface}, Schnittstelle zwischen Softwaresystemen) oder eine Datenbankoperation.
  \item \fachbegriff{Send/Receive Task} -- expliziter Nachrichtenaustausch über System- oder Pool-Grenzen.
  \item \fachbegriff{Business Rule Task} -- delegiert Entscheidung an eine Decision Engine (DMN).
  \item \fachbegriff{Timer Event} -- Reminder, Eskalation, Deadlines (``nach 3 Tagen erinnern, nach 5 Tagen eskalieren'').
  \item \fachbegriff{Message Event} -- asynchrone Auslöser, z.\,B.\ Eingang eines externen Events.
  \item \fachbegriff{Error Event} -- Fehlerbehandlung mit Boundary Events an User/Service Tasks.
  \item \fachbegriff{Datenobjekte und Prozess-Variablen} -- typisierte Datenstruktur, die zwischen Tasks weitergegeben wird.
\end{itemize}

\subsection{Gegenüberstellung}

\begin{tabularx}{\linewidth}{@{}l X X@{}}
\toprule
& \textbf{Business-BPMN} & \textbf{Execution-BPMN} \\
\midrule
\textbf{Adressat} & Fachseite, Management, Audit-Übersicht & Workflow-Engine, IT, Operation \\
\textbf{Tiefe} & Hauptpfad, wichtige Varianten & jeder Pfad vollständig durchführbar \\
\textbf{Gateways} & Benennung der Entscheidung & formale Bedingung in FEEL/Skript \\
\textbf{Tasks} & generisch (``Rechnung prüfen'') & typisiert (User/Service/Script/BR) \\
\textbf{Daten} & implizit oder annotiert & Variablen mit Typ, Datenobjekte \\
\textbf{Fehlerfälle} & ``geht zurück an Sachbearbeitung'' & Error Boundary Event, Eskalation \\
\textbf{Versionierung} & semantisch (v3.1 = Geschäftsänderung) & technisch (Deployment-Version) \\
\bottomrule
\end{tabularx}

\medskip

\noindent Beide Modelle haben Existenzrecht. Es ist Senior-Praxis, sie \emph{nicht} zu vermischen: das eine ist ein Kommunikationsobjekt, das andere ein Ausführungsobjekt. Die Brücke ist eine \fachbegriff{Übergabe-Routine}, die definiert, welche Teile des Business-BPMN in Execution-BPMN überführt werden -- und welche bewusst nicht.

\begin{merkbox}[Zwei Modelle, ein Owner]
Business- und Execution-BPMN sollten denselben \emph{Owner} haben. Wenn ``die Fachseite'' das Business-Modell pflegt und ``die IT'' das Execution-Modell, driften beide auseinander, sobald Änderungen kommen. Geteilte Verantwortung ist Geheimnis um Verantwortung.
\end{merkbox}

% --- 5. AUTOMATIONS-SLICE ------------------------------------------------
\section{Den automatisierbaren Slice finden}

Nicht jeder Prozess gehört in eine Engine. Die Kunst besteht darin, in einem dokumentierten Prozess die Teilstrecke zu identifizieren, die \emph{realistisch} und \emph{wertvoll} automatisierbar ist. Diese Teilstrecke nennen wir den \fachbegriff{Automations-Slice}. \quelle{vom Brocke \& Mendling, 2021}

\subsection{Sechs Kriterien für einen guten Slice}

\begin{enumerate}
  \item \textbf{Klarer Start und klares Ende.} Was triggert den Slice, wann ist er fertig? Wenn diese Grenzen verschwimmen, wird der Workflow zur Sammlung von Sonderlocken.
  \item \textbf{Eindeutige Daten.} Welche Felder kommen rein, welche werden zwischen Tasks weitergegeben, welche werden geschrieben? Ohne Datenklarheit kann keine Service Task funktionieren.
  \item \textbf{Bekannte Fehlerfälle.} Was passiert bei Timeout, fehlenden Daten, abgelehntem Antrag? Mindestens die häufigsten drei Fehlerpfade müssen vor Go-Live definiert sein.
  \item \textbf{Hohe Frequenz und niedrige Varianz.} Ein Workflow, der hundertmal pro Tag in fast immer derselben Form läuft, lohnt sich. Ein Workflow für drei Sonderfälle pro Jahr ist eine Checkliste.
  \item \textbf{Vorhandene Schnittstellen.} Wenn die Service Tasks auf nicht existierende APIs zugreifen sollen, ist die Automatisierung ein API-Projekt mit Workflow-Dekoration -- das muss man wissen.
  \item \textbf{Messbarer Nutzen.} Eine KPI vor Go-Live (Durchlaufzeit, Fehlerquote, Personenstunden) und ein realistischer Zielwert.
\end{enumerate}

\subsection{Beispiel: Rechnungsfreigabe als Slice}

Nehmen wir den klassischen Mini-Prozess ``Rechnungsfreigabe'':

\begin{enumerate}
  \item Rechnung kommt per E-Mail.
  \item Buchhaltung erfasst Grunddaten.
  \item Wenn Betrag $> 5.000\,\text{\euro}$, muss Teamleitung freigeben.
  \item Wenn Kostenstelle fehlt, geht es zurück an die Anforderungsstelle.
  \item Nach Freigabe wird im ERP (\emph{Enterprise Resource Planning}, unternehmensweites Ressourcenplanungssystem) gebucht und Zahlung angestoßen.
  \item Bei Fristüberschreitung wird erinnert und nach 3 Tagen eskaliert.
\end{enumerate}

\subsection{Slice-Analyse}

Wir markieren typisiert: \textit{User Tasks} (Menschenarbeit), \textit{Service Tasks} (Systemaktionen), \textit{Entscheidungen}, \textit{Fehlerfälle}, \textit{Fristen}.

\begin{itemize}
  \item Schritt 1 (Rechnung per E-Mail) -- \emph{Trigger} (Message Event aus Mailbox), evtl.\ Vorverarbeitung Service Task.
  \item Schritt 2 (Grunddaten erfassen) -- \emph{User Task} Buchhaltung, mit Pflichtfeldern (Betrag, Kostenstelle, Lieferant).
  \item Schritt 3 (Schwellwert-Entscheidung) -- \emph{Gateway} (Business Rule Task möglich).
  \item Schritt 3a (Freigabe Teamleitung) -- \emph{User Task} mit SLA.
  \item Schritt 4 (Kostenstelle fehlt) -- \emph{Error Path} mit Rückgabe an Anforderungsstelle.
  \item Schritt 5 (ERP-Buchung) -- \emph{Service Task} (API ERP).
  \item Schritt 5b (Zahlung) -- \emph{Service Task} (API Payment).
  \item Schritt 6 (Frist) -- \emph{Timer Boundary Event} an User Task, mit Reminder nach 24\,h und Eskalation nach 72\,h.
\end{itemize}

Datenminimum (drei Felder, die der Workflow zwingend braucht): \fachbegriff{Betrag}, \fachbegriff{Kostenstelle}, \fachbegriff{Freigeber-Rolle}.

\subsection{Ausnahmepolitik -- bewusst nicht automatisieren}

Es ist Senior-Praxis, im ersten Release nicht alle Ausnahmen abzubilden. Eine sinnvolle Ausnahmepolitik benennt explizit, welche Fälle \emph{organisatorisch} und nicht \emph{technisch} abgefangen werden. Beispiele für die Rechnungsfreigabe:

\begin{itemize}
  \item Lieferanten ohne Stammdatensatz: erst manuell in ERP anlegen, dann Workflow neu starten.
  \item Fremdwährungen: vorläufig nicht automatisiert, manueller Prozess.
  \item Sammelrechnungen mit mehreren Kostenstellen: vorläufig manuell, in Release 2 einplanen.
\end{itemize}

\begin{eselsbox}[Eselsbrücke: \akzent{S-D-F-F}]
Vor jedem Automations-Slice prüfen:\\[2pt]
\textbf{S}tart/Ende -- klar abgrenzbar?\\
\textbf{D}aten -- Minimum-Set definiert?\\
\textbf{F}ehlerfälle -- mindestens drei dokumentiert?\\
\textbf{F}requenz -- hoch genug, dass es sich lohnt?\\[2pt]
\emph{Wenn auch nur ein S-D-F-F-Buchstabe wackelt, ist es kein Slice, sondern ein Wunschzettel.}
\end{eselsbox}

% --- 6. KPIs UND STEUERUNG -----------------------------------------------
\section{KPIs für Repository und Workflows}

Was nicht gemessen wird, wird nicht gesteuert. Aber: was falsch gemessen wird, erzeugt Fehlanreize. KPIs für Repository und Execution sollten zugleich \emph{wenig}, \emph{wirksam} und \emph{nicht gameable} sein. \quelle{Reinkemeyer, 2020}

\subsection{KPI-Kandidaten für das Repository}

\begin{itemize}
  \item \fachbegriff{Approval-Quote.} Anteil der Modelle in Status ``Approved'' (Ziel: $>80\,\%$). Niedrig $\rightarrow$ Reviews stocken.
  \item \fachbegriff{Review-Aktualität.} Anteil der Modelle mit Review jünger als 12 Monate. Niedrig $\rightarrow$ Repository altert.
  \item \fachbegriff{Owner-Vollständigkeit.} Anteil der Modelle mit eindeutig benanntem Owner (Ziel: 100\,\%).
\end{itemize}

\subsection{KPI-Kandidaten für Workflow-Execution}

\begin{itemize}
  \item \fachbegriff{Durchlaufzeit (DLZ).} Zeit von Start- bis End-Event. Owner: Process Owner. Achtung: SLA-Definition muss vorher klar sein.
  \item \fachbegriff{Eskalationsquote.} Anteil der Cases, die mindestens einmal eskaliert wurden. Steigend $\rightarrow$ Slice zu eng geschnitten oder Ressourcen fehlen.
  \item \fachbegriff{Rework-Quote.} Anteil der Cases mit Korrekturschleife. Steigend $\rightarrow$ Datenqualität oder Pflichtfelder am Eingang prüfen.
  \item \fachbegriff{Auto-Resolution-Rate.} Anteil vollständig automatisch abgeschlossener Cases ohne manuelle Intervention.
\end{itemize}

\subsection{Anti-Gaming und Owner}

Jede KPI braucht einen Owner -- und ein Bewusstsein dafür, wie sie verfälscht werden könnte:

\begin{tabularx}{\linewidth}{@{}>{\bfseries}lXl@{}}
\toprule
KPI & Mögliches Gaming & Gegenmaßnahme \\
\midrule
DLZ & Cases früh ``vorläufig'' schließen & gekoppelt mit Rework-Quote \\
Eskalationsquote & Eskalationen umgehen oder umetikettieren & Audit-Trail-Check Stichprobe \\
Auto-Resolution & einfache Cases bevorzugen, schwere blockieren & Verteilung nach Komplexität \\
Approval-Quote & ``Quick Approvals'' ohne echte Reviews & Reviewer-Vier-Augen-Pflicht \\
\bottomrule
\end{tabularx}

\medskip

\begin{merkbox}[KPI braucht Kontext]
Eine KPI ohne Owner ist Statistik. Eine KPI ohne Gegenkontrolle ist Einladung zum Gaming. Eine KPI ohne Maßnahmen-Verknüpfung ist Dekoration. Senior-Steuerung verlangt: Owner, Gegenkontrolle, Maßnahmen-Verknüpfung -- für jede einzelne Kennzahl.
\end{merkbox}

% --- 7. FALLSTUDIE -------------------------------------------------------
\section{Fallstudie: MediSupply Services -- E2E-Repository und Pilot}

Wir betrachten ein realistisches Szenario, das die Themen des Tages verbindet -- Repository-Design, Pilot-Auswahl, Automations-Slice und KPI-Steuerung. \emph{MediSupply Services} ist ein B2B-Anbieter im Medizinprodukte-Umfeld mit mehreren Standorten. Die Probleme sind typisch: uneinheitliche Prozesse, schwankende Durchlaufzeiten, lückenhafte Audit-Nachweise. Ein Teilprozess soll schnell automatisiert werden, aber IT-Kapazität ist knapp.

\subsection{Ausgangslage und Restriktionen}

\begin{itemize}
  \item Zeit: 8 Wochen für Pilot.
  \item Budget: 90\,000\,\text{\euro}.
  \item ERP-Schnittstelle: existiert, aber nicht vollständig dokumentiert.
  \item Standorte: drei in DACH, jeweils mit eigenen Varianten.
\end{itemize}

\subsection{Schritt 1 -- E2E-Prozesslandkarte}

Wir entwerfen eine Landkarte mit fünf Kernprozessen:

\begin{enumerate}
  \item \emph{Lead-to-Order} -- Anfrage bis Auftragsbestätigung.
  \item \emph{Order-to-Delivery} -- Auftragsabwicklung bis Auslieferung.
  \item \emph{Invoice-to-Cash} -- Rechnungsstellung bis Zahlungseingang.
  \item \emph{Quality-to-CAPA} -- Qualitätsereignis bis Korrekturmaßnahme.
  \item \emph{Hire-to-Onboarding} -- Personalprozess.
\end{enumerate}

\subsection{Schritt 2 -- Pilot wählen}

Wir wählen \fachbegriff{Lead-to-Order} als Pilot. Begründung:

\begin{itemize}
  \item Hohe Frequenz (mehrere hundert Cases pro Monat).
  \item Klare Start-Ende-Definition (Anfrage rein, Bestätigung raus).
  \item Existierende Stammdaten (Kunde, Produkt) -- Datenminimum greifbar.
  \item Sichtbare DLZ-Probleme (Beschwerden aus Vertrieb).
\end{itemize}

Verworfene Alternative: \emph{Invoice-to-Cash}. Zwar hohe Wirkung auf Cashflow, aber abhängig von der unklaren ERP-Schnittstelle -- zu hohes Risiko im 8-Wochen-Zeitraum.

\subsection{Schritt 3 -- Repository-Struktur für den Pilot}

\begin{itemize}
  \item Ebene 1: ``Lead-to-Order'' (Owner: Head of Sales Ops, KPI-Set: DLZ, Konversionsrate, Reklamationsquote).
  \item Ebene 2 (Subprozesse): Anfrageannahme, Bonitäts-/Verfügbarkeitscheck, Angebotserstellung, Vertragsabschluss, Auftragsanlage.
  \item Ebene 3: SOP Anfrageannahme (Pflichtfelder), Eskalationsleitfaden Bonität, Angebots-Templates pro Produktkategorie.
\end{itemize}

Metadaten-Standard (fünf Pflichtfelder): Owner, Gültigkeit, Status, Standort, Risikoklasse. Zusätzlich: KPI-Verlinkung und letzte Änderung mit Approver-Name.

\subsection{Schritt 4 -- Automationskandidat identifizieren}

Innerhalb von Lead-to-Order ist der Subprozess \emph{Bonitäts- und Verfügbarkeitscheck} der heißeste Kandidat:

\begin{itemize}
  \item Start: Anfrage liegt vor, Kundendaten erfasst.
  \item Service Tasks: Bonitätsanfrage an externen Provider (sync), Verfügbarkeitscheck an ERP (sync).
  \item XOR: Bonität ok? Verfügbar?
  \item User Task bei Sonderfall: Sales-Mitarbeiter validiert manuell.
  \item Ende: Ergebnis-Datensatz für Angebotserstellung übergeben.
\end{itemize}

\subsection{Schritt 5 -- Execution-Skizze (Camunda-Style)}

\begin{quote}
\textbf{Start} (Message Event ``Anfrage geprüft'') \\
\textrightarrow{} \emph{Service Task ``Bonität abfragen''} (API extern, Timeout 5\,s, 2 Retries) \\
\textrightarrow{} \emph{Service Task ``Verfügbarkeit prüfen''} (API ERP, Idempotent, Case-ID) \\
\textrightarrow{} \textbf{XOR-Gateway} ``Beide ok?'' \\
\qquad \emph{Ja}: \textrightarrow{} \emph{Service Task ``Ergebnis speichern''} \textrightarrow{} \textbf{Ende} ``Check positiv'' \\
\qquad \emph{Nein}: \textrightarrow{} \emph{User Task ``Manuelle Prüfung''} (SLA 4\,h, Eskalation an Sales-Lead) \\
\qquad\quad \textrightarrow{} \emph{XOR ``Freigegeben?''} \textrightarrow{} \textbf{Ende} ``Check positiv'' oder ``Check negativ''.
\end{quote}

Fehlerfälle: API-Timeout (Retry, dann manuelle Prüfung), unbekannter Kunde (zurück an Sales mit Hinweis), ERP nicht verfügbar (Timer Event, Wiederholung nach 15\,min, dann Eskalation).

\subsection{Schritt 6 -- KPIs für den Pilot}

\begin{itemize}
  \item \fachbegriff{DLZ Bonitäts-/Verfügbarkeitscheck} -- Ziel: $<30\,\text{min}$ in 90\,\% der Fälle. Owner: Process Owner Lead-to-Order.
  \item \fachbegriff{Eskalationsquote manuelle Prüfung} -- Ziel: $<15\,\%$. Owner: Sales-Lead.
  \item \fachbegriff{Rework-Quote} -- Cases, die nach Check zurück an Sales gehen. Ziel: $<5\,\%$. Owner: Sales Ops.
\end{itemize}

\subsection{Lessons aus der Fallstudie}

\begin{itemize}
  \item Pilotwahl ist kein Lostopf, sondern eine Bewertung nach Frequenz, Klarheit, Datenverfügbarkeit, Schnittstellen-Reife.
  \item Variantenpolitik wird früh, nicht spät, entschieden -- sonst wachsen drei Modelle parallel.
  \item Automation ohne KPI ist Selbstzweck. KPI ohne Owner ist Statistik. Owner ohne Maßnahmenrecht ist Theater.
\end{itemize}

% --- 8. MANAGEMENT-PERSPEKTIVE -------------------------------------------
\section{Management-Perspektive: Repository als Betriebssystem, Execution als Vertrag}

Auf Wissens-Ebene sind Repository und Execution Notations- und Werkzeug-Themen. Auf Anwendungs-Ebene sind sie Handwerk. Auf Management-Ebene sind sie \emph{Operating-Model-Entscheidungen} -- mit Auswirkungen auf Verantwortlichkeiten, Kosten, Risiken und Kultur.

\subsection{Zwei zentrale Zielkonflikte}

\begin{enumerate}
  \item \textbf{Governance-Strenge vs.\ Akzeptanz.} Strenge Reviews und Freigabeprozesse erzeugen Qualität -- und Widerstand. Zu lockere Regeln erzeugen Wildwuchs. Praxis: Minimum-Set definieren, alles darüber bewusst entscheiden -- nicht ``alles oder nichts''.
  \item \textbf{Standardisierung vs.\ Standortautonomie.} Wenn jeder Standort seine Variante behalten will, gibt es keine Vergleichbarkeit. Wenn ein zentrales Modell alle Standorte zwingt, geht lokale Klugheit verloren. Praxis: \emph{Guardrails statt Wildwuchs} -- harte Regeln für 5 nicht verhandelbare Standards, lokale Freiheit innerhalb dieser Leitplanken.
\end{enumerate}

\subsection{Entscheidungen unter Unsicherheit}

In Repository- und Workflow-Initiativen sind die Entscheidungen mit langer Halbwertszeit besonders teuer:

\begin{itemize}
  \item \fachbegriff{Tool-Stack.} Einmal etabliert, kostet ein Tool-Wechsel mindestens ein Jahr. Vorher klären: braucht es ein BPMS, oder reicht eine Modellierungs-Plattform mit Workflow-Anschluss?
  \item \fachbegriff{Owner-Modell.} Wer ist End-to-End-Owner -- ein Fachbereich, ein Querschnittsamt, oder eine Linien-Rolle? Diese Entscheidung prägt die nächsten drei bis fünf Jahre.
  \item \fachbegriff{Variantenpolitik.} ``Ein Modell für alle'' vs.\ ``ein Modell pro Standort'' -- jede Antwort hat sechsstellige Folgekosten in Pflege.
\end{itemize}

Für jede dieser Entscheidungen gilt: schriftlich treffen, mit verworfener Alternative dokumentieren, mit Frühwarnsignal versehen (``wenn Bedingung X eintritt, prüfen wir die Entscheidung erneut'').

\subsection{Risikoarchitektur eines Workflow-Programms}

\begin{enumerate}
  \item \textbf{Daten-Risiko.} Workflow scheitert an Datenqualität. Mitigation: Pflichtfelder im Eingangs-Formular, Service-Task ``Datenvalidierung'' vor jedem Service Task mit externem System.
  \item \textbf{Schnittstellen-Risiko.} APIs ändern sich, ohne dass es jemand merkt. Mitigation: Versionierung, Contract Tests (konzeptionell), API-Owner mit Change-Notification-Pflicht.
  \item \textbf{Compliance-Risiko.} Audit-Findings durch fehlende Trails. Mitigation: jedes Approved-Modell mit Datum/Approver, jeder Workflow-Lauf mit unveränderlichem Audit-Log.
  \item \textbf{Akzeptanz-Risiko.} Operative Teams umgehen den Workflow durch Schatten-Prozesse. Mitigation: frühzeitige Einbindung, ``Was wird für mich besser?''-Argumentation, klare Eskalationswege.
\end{enumerate}

\subsection{Senior-Shift: vom Tool-Verstehen zur Verantwortungsarchitektur}

Der Wechsel von ``ich modelliere ein Diagramm'' zu ``ich verantworte ein Steuerungssystem'' erkennt man an drei Kennzeichen:

\begin{itemize}
  \item Man entscheidet die Variantenpolitik \emph{vor} dem ersten Tool-Klick.
  \item Man kennt die zwei Governance-Regeln, die 80\,\% des Chaos verhindern -- und verteidigt sie auch unter Druck.
  \item Man kann in 90 Sekunden gegenüber einem Vorstand begründen, warum genau dieser Slice automatisiert wurde und nicht ein anderer.
\end{itemize}

\begin{merkbox}[Schluss-Merksatz für Tag 4]
Ein Repository ohne Owner ist eine Ablage. Eine Execution ohne Audit Trail ist eine Black Box. Ein Workflow ohne KPI ist eine Verlagerung des Problems. Senior-Praxis verbindet alle drei: \emph{Owner steuert das Repository, Audit Trail sichert die Execution, KPI macht beides messbar.}
\end{merkbox}

% --- 9. TAKE-AWAYS -------------------------------------------------------
\section{Take-aways aus Tag 4}

\begin{enumerate}
  \item \textbf{Repository = Betriebssystem.} Drei Ebenen (E2E \textrightarrow{} Subprozess \textrightarrow{} Arbeitsdokument), fünf Pflicht-Metadaten, vier Rollen -- mehr braucht es zum Start nicht.
  \item \textbf{Execution = Vertrag.} Daten, Fehler, SLAs, Eskalation -- jeder Punkt schriftlich, sonst trägt im Ernstfall niemand.
  \item \textbf{Business-BPMN $\neq$ Execution-BPMN.} Zwei Modelle, ein Owner, eine definierte Übergabe-Routine.
  \item \textbf{Slice statt All-in.} Ein klar abgegrenzter Automations-Slice mit messbarem Nutzen schlägt jede ``großen Wurf''-Ankündigung.
  \item \textbf{Zwei Governance-Regeln verhindern 80\,\% Chaos.} Owner + Approved-Status mit Datum sowie Change Requests mit Reviewer.
  \item \textbf{Variantenpolitik früh entscheiden.} Attribute oder separate Modelle -- die Wahl prägt Pflegekosten über Jahre.
  \item \textbf{KPI braucht Owner + Anti-Gaming.} Jede Kennzahl mit verantwortlicher Person und Gegenkontrolle, sonst entsteht Reporting ohne Steuerung.
  \item \textbf{Guardrails statt Wildwuchs.} Zentrale Standards für das Minimum, lokale Freiheit innerhalb -- Standortautonomie ohne Beliebigkeit.
\end{enumerate}

% --- 10. LITERATUR -------------------------------------------------------
\clearpage
\section{Literatur}

Im Skript verwendete und für die Vertiefung empfohlene Quellen. Zitierstil: APA-7.

\begin{description}[style=nextline,leftmargin=18pt,labelindent=0pt,itemsep=4pt]
  \item[Allweyer, T. (2020).] \emph{BPMN 2.0 -- Business Process Model and Notation: Einführung in den Standard für die Geschäftsprozessmodellierung} (4.\ Aufl.). BoD.
  \item[Dumas, M., La Rosa, M., Mendling, J., \& Reijers, H.\,A. (2018).] \emph{Fundamentals of business process management} (2nd ed.). Springer. \href{https://doi.org/10.1007/978-3-662-56509-4}{https://doi.org/10.1007/978-3-662-56509-4}
  \item[Freund, J., \& Rücker, B. (2019).] \emph{Real-life BPMN: Includes an introduction to DMN} (4th ed.). CreateSpace.
  \item[Harmon, P. (2019).] \emph{Business process change: A business process management guide for managers and process professionals} (4th ed.). Morgan Kaufmann.
  \item[Object Management Group. (2014).] \emph{Business Process Model and Notation (BPMN), Version 2.0.2} (formal/2013-12-09). \href{https://www.omg.org/spec/BPMN/2.0.2/}{https://www.omg.org/spec/BPMN/2.0.2/}
  \item[Reinkemeyer, L. (Hrsg.). (2020).] \emph{Process mining in action: Principles, use cases and outlook}. Springer. \href{https://doi.org/10.1007/978-3-030-40172-6}{https://doi.org/10.1007/978-3-030-40172-6}
  \item[Rosemann, M., \& vom Brocke, J. (2010).] The six core elements of business process management. In J. vom Brocke \& M. Rosemann (Hrsg.), \emph{Handbook on business process management 1} (S.~107--122). Springer. \href{https://doi.org/10.1007/978-3-642-00416-2_5}{https://doi.org/10.1007/978-3-642-00416-2\_5}
  \item[Scheer, A.-W. (2002).] \emph{ARIS -- Vom Geschäftsprozess zum Anwendungssystem} (4.\ Aufl.). Springer.
  \item[vom Brocke, J., \& Mendling, J. (Hrsg.). (2021).] \emph{Business process management cases: Digital innovation and business transformation in practice} (Vol.~2). Springer. \href{https://doi.org/10.1007/978-3-662-63047-1}{https://doi.org/10.1007/978-3-662-63047-1}
\end{description}

% --- 11. GLOSSAR ---------------------------------------------------------
\clearpage
\section{Glossar}

Die wichtigsten Begriffe des Tages -- kurz definiert, in alphabetischer Reihenfolge.

\begin{description}[style=nextline,leftmargin=0pt,labelindent=0pt,itemsep=4pt]
  \item[Approval-Quote] Anteil der Modelle in Status ``Approved''. Indikator für Governance-Aktivität.
  \item[Approver] Rolle, die finale Modellversionen freigibt. Häufig der Process Owner.
  \item[Audit Trail] Unveränderliches Log aller Schritte und Entscheidungen eines Workflow-Laufs. Pflicht in regulierten Umfeldern.
  \item[Automations-Slice] Klar abgegrenzte Teilstrecke eines Prozesses, die als Workflow ausgeführt wird. Erfordert Start/Ende, Datenminimum, Fehlerfälle.
  \item[BPMS] \emph{Business Process Management System.} Kombiniert Modellierung, Ausführung, Überwachung und Optimierung in einer Plattform.
  \item[Business-BPMN] BPMN-Variante für Fach- und Management-Sicht. Hauptpfade, Lanes, Annotations -- nicht zwingend ausführbar.
  \item[Camunda] Verbreitete Workflow-Engine, die BPMN 2.0 ausführt. Open-Core mit Enterprise-Erweiterungen.
  \item[Change Request] Strukturierter Antrag zur Änderung eines Modells, mit Reviewer und Freigabe.
  \item[Datenminimum] Kleinste Menge an Feldern, die ein Workflow zwingend braucht, um ausgeführt zu werden.
  \item[Definition of Done (DoD)] Liste prüfbarer Kriterien, ab denen ein Modell oder Workflow als fertig gilt.
  \item[E2E-Prozess] End-to-End-Prozess: vollständiger Wertschöpfungsstrang von Trigger bis Outcome.
  \item[Eskalationsquote] Anteil der Cases, die mindestens einmal eskaliert wurden. KPI für Workflow-Steuerung.
  \item[Execution-BPMN] BPMN-Variante für Workflow-Engine. Jeder Pfad vollständig, jede Bedingung formal, Tasks typisiert.
  \item[Freigabeprozess] Lifecycle eines Modells: Draft \textrightarrow{} Review \textrightarrow{} Approved \textrightarrow{} Retired.
  \item[Governance] Regeln, Rollen und Prozesse, die Modelle steuerbar und auditierbar machen.
  \item[Guardrails] Verbindliche Leitplanken, innerhalb derer lokale Autonomie zulässig ist.
  \item[MVP] \emph{Minimum Viable Product.} Erste Version, die Lerneffekt erzeugt -- nicht alle Features.
  \item[Owner (Process Owner)] Person mit fachlicher End-to-End-Verantwortung für einen Prozess. Pflicht-Metadatum jedes Modells.
  \item[Repository] Verwaltete Sammlung von Prozessmodellen und Artefakten mit Metadaten, Versionen und Lebenszyklus.
  \item[Service Task] BPMN-Element: ein Schritt, der durch ein System ausgeführt wird, typischerweise API-Aufruf.
  \item[SLA] \emph{Service Level Agreement.} Vereinbarung über Antwort- oder Durchlaufzeiten.
  \item[Status] Zustand eines Modells im Lebenszyklus: Draft, Review, Approved, Retired.
  \item[Timer Event] BPMN-Element für zeitgesteuerte Aktionen: Reminder, Eskalation, Deadlines.
  \item[User Task] BPMN-Element: ein Schritt, der durch einen Menschen ausgeführt wird, typischerweise mit Form.
  \item[Variantenattribut] Mechanismus, um Modellvarianten als Datenfeld am gemeinsamen Modell zu führen.
\end{description}

\vspace{1cm}
\noindent{\color{lpAccent}\rule{\linewidth}{0.6pt}}\\[4pt]
{\footnotesize\sffamily\color{lpGray}Lernplatt \textperiodcentered{} by Can Siebert \textperiodcentered{} Kurs 71 (Dig 42) \textperiodcentered{} Tag 4 \textperiodcentered{} \textcopyright{} 2026}

\end{document}
